\section{奖励机制：内生 reward 与 Multi-Agent}

本章进入 reward。姚顺雨提到，Agent 发展的关键方向之一，是让它拥有自己的 reward，能自己探索；另一个方向是 Multi-Agent，让它们之间形成组织结构。这里的 reward 不只是 RL 公式，而是 Agent 能否从环境中获得有效反馈、能否不只模仿人类数据、能否在任务中自我改进的核心。

\begin{figure}[H]
\centering
\includegraphics[width=0.94\textwidth]{reward-multi-agent.png}
\caption{Agent 两个关键方向：拥有自己的 reward，与多智能体形成组织结构。自制概念图，依据 00:27:19--01:03:00 对谈内容整理。}
\end{figure}

\begin{knowledgebox}{读图：reward 和 multi-agent 是两条不同扩展路径}
Reward 让单个 Agent 能探索、试错和优化；Multi-Agent 让多个 Agent 通过分工、协作和竞争形成更复杂的组织。前者解决学习信号，后者解决组织结构。
\end{knowledgebox}

\subsection{Reward 难在哪里}

上一节把 reward 和 multi-agent 分成两条扩展路径，本节先看 reward 为什么难。现实任务的 reward 很难定义。棋类游戏有明确胜负，代码任务有测试，网页任务可能有完成状态；但很多现实任务的目标长期、稀疏、主观或不可完全观测。reward 设计不好，Agent 可能优化错目标，或者钻评价漏洞。

\begin{figure}[H]
\centering
\includegraphics[width=0.96\textwidth]{reward-definition-problem.png}
\caption{奖励定义问题：现实任务的 reward 难在长期、稀疏和不可完全观测。自制概念图，依据 00:40:00--01:03:00 对谈内容整理。}
\end{figure}

\begin{warningbox}{Reward hacking 的风险}
如果 reward 只奖励表面结果，Agent 可能学会投机。例如只奖励点击率，它可能制造标题党；只奖励任务完成，它可能绕过权限；只奖励测试通过，它可能写脆弱代码。企业和真实世界任务尤其需要审计和安全约束。
\end{warningbox}

\begin{importantbox}{Reward 的直觉公式}
\[
\text{learning signal} = r(s_t, a_t, s_{t+1})
\]
其中，\(s_t\) 是行动前状态，\(a_t\) 是 Agent 的动作，\(s_{t+1}\) 是行动后状态，\(r\) 是评价这个变化的奖励函数。现实任务难在：状态不完整、动作后果延迟、评价标准常常由人和组织共同决定。
\end{importantbox}

\subsection{Multi-Agent：从个体能力到组织结构}

Multi-Agent 的关键不是“多开几个模型实例”，而是让 Agent 之间形成分工、通信、协作、竞争和组织结构。它接近人类组织，也会带来新的问题：谁分配任务，谁验证结果，谁承担冲突，如何避免群体幻觉。

\begin{knowledgebox}{术语消化：Reward 与 Multi-Agent}
\begin{tabularx}{\textwidth}{p{0.22\textwidth}p{0.32\textwidth}X}
\toprule
概念 & 解决的问题 & 常见风险 \\
\midrule
External reward & 环境或人类给出的外部评价 & 昂贵、稀疏、容易被过拟合。 \\
Intrinsic reward & Agent 自己形成探索信号 & 可能偏离人类目标。 \\
Multi-Agent & 多个 Agent 分工协作 & 通信成本、冲突和责任边界。 \\
Verification & 验证 Agent 输出是否可靠 & 验证器本身也可能出错。 \\
\bottomrule
\end{tabularx}
\end{knowledgebox}

Reward 和 multi-agent 的难点还在于它们会改变系统的“社会性”。单个 Agent 犯错时，问题可以追溯到单个策略；多个 Agent 协作时，错误可能来自任务分配、通信误解、局部最优、验证器缺陷或激励不一致。越接近真实组织，Agent 系统越需要治理层，而不是只需要更强模型。

\begin{warningbox}{多智能体不是自动等于更强}
多 Agent 系统可能带来分工和并行，也可能带来更高通信成本、更复杂失败模式和责任不清。只有当任务天然可分解、验证机制清楚、协作协议稳定时，多智能体才更可能超过单 Agent。
\end{warningbox}

\subsection{本章小结}

Reward 决定 Agent 学什么，Multi-Agent 决定 Agent 如何组织。下半场的难点不是让模型更会聊天，而是让系统能在真实环境里安全地探索和协作。

